\documentclass[10pt,letterpaper]{article}
\usepackage{cornell-notes}

\title{Clause 29.11.02.03: Time Zone Database Access}
\author{}
\date{}
\setDocTitle{Clause 29.11.02.03: Time Zone Database Access}
\setDocSubtitle{C++ 2024}
\setDocOwner{C++ 2024 Cornell Notes}
\setDocDate{\today}
\setCornellCollection{C++ 2024 Cornell Notes}
\setCornellUnitType{Clause}
\setCornellUnitNumber{29.11.02.03}
\setCornellUnitTitle{Time Zone Database Access}
\hypersetup{pdftitle={Clause 29.11.02.03: Time Zone Database Access},pdfsubject={Cornell notes},pdfauthor={}}

\begin{document}
\maketitle
\begin{CornellOverviewBox}{Overview}
\textbf{Classification:} Standard library - Normative\\[2pt]
\textbf{Learning objective:} Understand the rules, terminology, and semantic consequences associated with Time zone database access.
\end{CornellOverviewBox}\begin{CornellSummaryBox}{Summary}
{\sffamily\bfseries\color{Accent} Summary}\par\vspace{3pt}Section 29.11.2.3 focuses on Time zone database access. Apply the specified interface together with its mandates, effects, result conditions, exception behavior, and complexity constraints. When applying it, preserve the distinction between mandatory requirements, recommendations, permissions, and behavior deliberately left undefined, unspecified, or implementation-defined.
\end{CornellSummaryBox}

\end{document}
